%%%%%%%%%%%%%%%%%%%%%%%%%%%%%%%%%%%%%%%%%%%%%%%%%%%%%%%%%%%%%%%%%%%%%%%%%%%
% Manufactured FSI solution: construction, boundary data and checks.
% LaTeX-ready text for the verification paper (subsection + appendix).
% Notation of sectionMathModel.tex: v_F, D_S, sigma_S, P_S, p_F (kinematic),
% nu = mu_F/rho_F, OpenFOAM gradient convention (grad a)_ij = d a_j/d x_i,
% F = I + (grad_0 D_S)^T. Uses \bb (= \boldsymbol) from article.tex.
% Source of every formula: derive/mms_fsi.py (sympy/symengine), with the
% checks in derive/generated*.derive.log and derive/generated*.check.log.
%%%%%%%%%%%%%%%%%%%%%%%%%%%%%%%%%%%%%%%%%%%%%%%%%%%%%%%%%%%%%%%%%%%%%%%%%%%

\subsection{A manufactured solution for ALE fluid--structure interaction}
\label{sec:mmsc}

The manufactured solution couples a two-dimensional channel flow to a compliant, finite-strain wall.
It is closed form, starts from rest, and satisfies both interface conditions~\eqref{eqn:kinematic_interface} exactly, so that the coupling is tested without an interface source term.
The fluid occupies $0 < x < L$, $0 < y < h(x,t)$, with an inlet at $x = 0$, an outlet at $x = L$, a no-slip wall at $y = 0$ and the interface $y = h(x,t)$.
The solid occupies $0 < X < L$, $H < Y < H + T$ in its initial configuration, with the interface at $Y = H$, a clamped top at $Y = H + T$, and ends at $X = 0$ and $X = L$.
The parameters are $L = 2$, $H = 1$, $T = 1/2$; $\rho_F = 1$, $\mu_F = 1/100$; $\rho_{S,0} = 1$, $E_S = 10$, $\nu_S = 3/10$; and the period is $1$, $\omega = 2\pi$.
The density ratio is one, so added-mass effects are strong.
The motion is built from the time factors
\begin{equation}
  r(t) = \tfrac{1}{2}\left(1 - \cos\omega t\right), \qquad r_2(t) = \sin^2\omega t,
  \label{eq:mmsc_r}
\end{equation}
which vanish with their first derivatives at $t = 0$.
The displacement, the velocities and the pressure are therefore zero initially (the sources are not), and the calculation starts from rest on an undeformed mesh.

\paragraph{Interface motion, built from the solid side.}
The kinematic condition holds at the current position of each material interface point, $\bb{x}^\Gamma(X,t) = (x^\Gamma, y^\Gamma)$, whereas the fluid is written in Eulerian coordinates.
The construction therefore starts from the solid, and uses an interface parametrisation $X \mapsto x^\Gamma$ with a closed-form inverse.
The tangential motion is a M\"obius map of $[0,L]$ onto itself,
\begin{equation}
  x^\Gamma(X,t) = \frac{L\,q}{q + \kappa(1-q)}, \qquad
  \Xi(x,t) = \frac{L\,\kappa\,\hat{x}}{1 - \hat{x} + \kappa\,\hat{x}}, \qquad
  q = \frac{X}{L}, \quad \hat{x} = \frac{x}{L}, \quad \kappa(t) = 1 + \tfrac{1}{10}\,r_2(t),
  \label{eq:mmsc_mobius}
\end{equation}
where $\Xi$ is the inverse, $\Xi(x^\Gamma(X,t),t) = X$.
It fixes both end points and stretches the interface tangentially by up to $10\%$.
The normal motion is prescribed in Eulerian form, as the channel height
\begin{equation}
  h(x,t) = H + \hat\eta(x,t), \qquad
  \hat\eta = \tfrac{1}{10}\,\sin^2\!\left(\frac{\pi x}{L}\right) r(t)\left[1 + \tfrac{1}{2}\sin\!\left(\frac{2\pi x}{L} - \omega t\right)\right],
  \label{eq:mmsc_h}
\end{equation}
a pulsing wave with a travelling part.
The interface displacement is $\xi = x^\Gamma - X$ tangentially and $\eta(X,t) = \hat\eta(x^\Gamma(X,t),t)$ normally.
Prescribing $h$ in Eulerian rather than Lagrangian form keeps the swept area $\int h\,\mathrm{d}x$, which the fluid flux needs, elementary.
The factor $\sin^2(\pi x/L)$ makes $\hat\eta$ and its slope vanish at the ends, where the solid is clamped.
In the \emph{flat} variant the interface stays flat, $\hat\eta = \tfrac{1}{10}\,r(t)\left[1 + \tfrac12\sin\omega t\right]$ and $\kappa = 1$, so that the fluid mesh stays orthogonal at the interface; the solid ends then carry the exact displacement instead of a clamp, and the other definitions are unchanged.

\paragraph{Solid displacement and the shear identity.}
With $\hat{y} = (Y - H)/T$ and $g = (1 - \hat{y})(1 + \tfrac14\hat{y})$,
\begin{equation}
  D_{S,x} = \xi\,g + (Y - H)(1 - \hat{y})^2\,(m_0 + m_1), \qquad
  D_{S,y} = \eta\,g .
  \label{eq:mmsc_DS}
\end{equation}
$\bb{D}_S$ vanishes on the top and, in the full variant, on both ends, so that the interface is then the only solid boundary with non-zero data.
The stress is the compressible neo-Hookean law of solids4foam (\texttt{neoHookeanElastic}) in plane strain ($F_{zz} = 1$),
\begin{equation}
  \bb{\sigma}_S = \frac{\mu_S}{J}\,\mathrm{dev}\!\left(J^{-2/3}\bb{F}\bb{F}^T\right) + \frac{K_S}{2}\,\frac{J^2 - 1}{J}\,\bb{I},
  \qquad K_S = \lambda_S + \tfrac{2}{3}\mu_S ,
  \label{eq:mmsc_law}
\end{equation}
with $J = \det\bb{F}$, the three-dimensional deviator $\mathrm{dev}\,\bb{A} = \bb{A} - \tfrac13(\mathrm{tr}\,\bb{A})\bb{I}$, and the Lam\'e parameters $\mu_S = E_S/(2(1+\nu_S))$ and $\lambda_S = \nu_S E_S/((1+\nu_S)(1-2\nu_S))$.
For an arbitrary displacement the interface shear traction is of order $\mu_S$ times the strain; balancing it with $\mu_F = 10^{-2}$ would need unphysically large fluid shear rates.
For the law~\eqref{eq:mmsc_law}, the volumetric and the isotropic terms give no shear, and the deformed normal is parallel to $\bb{F}^{-T}\bb{N}$, with $|\bb{F}^{-T}\bb{N}| = |\bb{a}|/J$ in two dimensions.
The shear traction on the interface is therefore
\begin{equation}
  \bb{\tau}\cdot\bb{\sigma}_S\cdot\hat{\bb{n}} = \pm\,\mu_S\,J^{-2/3}\,\frac{\bb{a}\cdot\bb{b}}{|\bb{a}|^2},
  \qquad
  \bb{a} = \frac{\partial\bb{x}}{\partial X}, \quad \bb{b} = \frac{\partial\bb{x}}{\partial Y},
  \label{eq:mmsc_shear}
\end{equation}
where $\bb{\tau} = \bb{a}/|\bb{a}|$ is the unit tangent and $\hat{\bb{n}} = \hat{\bb{n}}_F$ (with the solid normal $\hat{\bb{n}}_S$ the sign is reversed).
The shear vanishes when the deformed material fibre $\bb{b}$ is orthogonal to the deformed interface.
On $Y = H$, $\bb{a} = (1 + \xi_{,X},\ \eta_{,X})$ and $\bb{b} = (\xi g_0' + m_0 + m_1,\ 1 + \eta g_0')$, with $g_0' = g'(H) = -3/(4T)$.
The choice
\begin{equation}
  m_0 = -\frac{\eta_{,X}\left(1 + \eta\,g_0'\right)}{1 + \xi_{,X}} - \xi\,g_0', \qquad
  m_1 = \tfrac{1}{50}\,\sin\!\left(\frac{2\pi X}{L}\right) r(t)
  \label{eq:mmsc_m}
\end{equation}
makes $\bb{a}\cdot\bb{b} = 0$ for $m_1 = 0$.
$m_1$ then adds a small, smooth shear of known size: at most $0.08$ in magnitude, against a normal traction of up to $3.0$ (maxima over a $401\times401$ grid in $X$ and $t$).
Over a sampling grid the solid reaches a Green--Lagrange strain (Frobenius norm) of $0.26$, a rotation of $0.23$\,rad and $0.69 \le J \le 1.09$.

\paragraph{Divergence-free fluid velocity.}
The fluid velocity is the curl of a stream function in the stretched coordinate $s = y/h(x,t)$,
\begin{equation}
  \bb{v}_F = \left(\frac{\partial\psi}{\partial y},\, -\frac{\partial\psi}{\partial x}\right), \qquad
  \psi = Q\,B(s) + h\,V_x\,C(s) + s^2(1-s)^2\,\Phi(x,t) + s^2(1-s)^3\,\Omega(x,s,t),
  \label{eq:mmsc_psi}
\end{equation}
so that $\bb{\nabla}\cdot\bb{v}_F = 0$ identically and no mass source is needed.
Here $\bb{V}(x,t) = \partial\bb{D}_S/\partial t\,(\Xi(x,t), H, t)$ is the interface velocity at the Eulerian abscissa, and $B = 3s^2 - 2s^3$ and $C = s^3 - s^2$ are Hermite functions.
Every term vanishes with its $s$-derivative at $s = 0$ (no slip).
At $s = 1$, $\psi = Q$ and $\partial\psi/\partial y = V_x$, and $v_{F,y} = -Q_{,x} + h_{,x}V_x$.
With $Q_{,x} = -h_{,t}$ this is $h_{,t} + h_{,x}V_x = V_y$, the rate of change of the height following the sliding interface point, so $\bb{v}_F = \bb{V}$ on the interface.
The flux
\begin{equation}
  Q(x,t) = r(t) + \frac{\partial}{\partial t}\int_x^L \hat\eta(x',t)\,\mathrm{d}x'
  \label{eq:mmsc_Q}
\end{equation}
prescribes a non-negative outlet flux $Q(L,t) = r(t)$; the inlet absorbs the change of volume of the channel, and the outlet never sees backflow.
The last term of~\eqref{eq:mmsc_psi} is a vortex train travelling downstream,
\begin{equation}
  \Omega = 8\,r(t)\,\sin^2\!\left(\frac{\pi x}{L}\right)\sin\!\left(\frac{2\pi x}{L} - \omega t\right)\left(1 + \tfrac{1}{2}\cos\pi s\right).
\end{equation}
It is faded out at both ends, and its weight $s^2(1-s)^3$ vanishes with two derivatives at $s = 1$, so it changes neither the interface velocity nor the interface traction.
The peak speed is about $2.1$, and the Reynolds number $\rho_F |\bb{v}_F|_{\max} H/\mu_F$ is about $200$.

\paragraph{Traction continuity without an interface source.}
On $y = h$, the unit normal and tangent are $\hat{\bb{n}}_F = (-h_{,x}, 1)/\sqrt{1 + h_{,x}^2}$ and $\bb{\tau} = (1, h_{,x})/\sqrt{1 + h_{,x}^2}$.
$\hat{\bb{n}}_F = -\hat{\bb{n}}_S$, where $\hat{\bb{n}}_S = \bb{F}^{-T}\bb{N}/|\bb{F}^{-T}\bb{N}|$ is the outward normal of the deformed solid and $\bb{N} = (0,-1)$.
The fluid traction $\rho_F\,\hat{\bb{n}}_F\cdot\bb{\sigma}_F^*$ is linear in two functions that are still free on the interface: $\Phi$ and the interface pressure.
The term $s^2(1-s)^2\Phi$ has zero value and slope at $s = 1$, so $\Phi$ leaves the interface velocity unchanged and enters the interface shear rate only through its second $s$-derivative.
The tangential balance is therefore algebraic in $\Phi$,
\begin{equation}
  c_0 + c_1\,\Phi = \bb{\tau}\cdot\bb{\sigma}_S\cdot\hat{\bb{n}}_F,
  \qquad c_1 = \frac{2\mu_F\left(1 + h_{,x}^2\right)}{h^2},
  \label{eq:mmsc_Phi}
\end{equation}
where $c_0$ is the fluid shear traction of the other terms of~\eqref{eq:mmsc_psi}, and $\bb{\sigma}_S$ is evaluated at $(\Xi(x,t), H, t)$.
This gives $\Phi$ in closed form.
The normal balance gives the interface pressure, and a term that vanishes on the interface completes the pressure field:
\begin{equation}
  \rho_F\,p_F = \rho_F\,p_\Gamma + (1 - s)\,\tfrac{1}{2}\,r(t)\cos\!\left(\frac{\pi x}{L}\right), \qquad
  \rho_F\,p_\Gamma = 2\mu_F\,\hat{\bb{n}}_F\cdot\bb{d}\cdot\hat{\bb{n}}_F - \hat{\bb{n}}_F\cdot\bb{\sigma}_S\cdot\hat{\bb{n}}_F,
  \label{eq:mmsc_p}
\end{equation}
with $\bb{d}$ the rate of strain of $\bb{v}_F$.
Both conditions~\eqref{eqn:kinematic_interface} then hold exactly.

\paragraph{Sources.}
The body forces are the residuals of the exact fields in Equations~\eqref{eqn:fluid_momentum_ale} and~\eqref{eqn:solid_momentum_TL}:
\begin{equation}
  \bb{b}_F = \frac{\partial\bb{v}_F}{\partial t} + \bb{v}_F\cdot\bb{\nabla}\bb{v}_F - \bb{\nabla}\cdot\bb{\sigma}_F^*, \qquad
  \rho_{S,0}\,\bb{b}_S = \rho_{S,0}\,\frac{\partial^2\bb{D}_S}{\partial t^2} - \bb{\nabla}_0\cdot\bb{P}_S^T .
  \label{eq:mmsc_sources}
\end{equation}
$\bb{b}_F$ is Eulerian, and applies unchanged on the moving mesh, since $\partial_t\bb{v}_F|_{\chi} + (\bb{v}_F - \bb{v}^\Gamma)\cdot\bb{\nabla}\bb{v}_F = \partial_t\bb{v}_F|_{\bb{x}} + \bb{v}_F\cdot\bb{\nabla}\bb{v}_F$ for any velocity $\bb{v}^\Gamma$ of the mesh points, with $\chi$ the mesh coordinates.
It is evaluated at the current cell centres at the new time level, and $\bb{b}_S$ at the initial cell centres.
The solid source per unit initial volume, $\rho_{S,0}\bb{b}_S$, is what the code applies.
Both are of the size of the physical terms: $|\bb{b}_F| \approx 10$ against $|\partial\bb{v}_F/\partial t| \approx 9$ and $|\bb{\nabla}p_F| \approx 8$ (maxima over a sampling grid, full variant).

\paragraph{Boundary data.}
Every prescribed boundary value and gradient is the exact one (Table~\ref{tab:mmsc_bc}); the interface data, and the momentum-balance and flux-consistent pressure gradients, are computed by the solver from the numerical fields.

\begin{table}[htbp]
  \centering
  \small
  \caption{Boundary conditions of the manufactured FSI solution. ``Exact'' means the manufactured value or normal gradient evaluated at the face centre at the new time level.}
  \label{tab:mmsc_bc}
  \begin{tabular}{lll}
    \toprule
    Boundary & Velocity / displacement & Pressure \\
    \midrule
    Inlet, $x = 0$ & exact $\bb{v}_F$ & QM: exact $\partial p_F/\partial n$; partitioned: flux-consistent gradient \\
    Outlet, $x = L$ & exact $\partial\bb{v}_F/\partial n$ & exact $p_F$ (sets the level) \\
    Wall, $y = 0$ & $\bb{v}_F = \bb{0}$ & QM: momentum balance; partitioned: flux-consistent gradient \\
    Interface & from the coupling & from the coupling \\
    Solid top & $\bb{D}_S = \bb{0}$ & --- \\
    Solid ends & full: $\bb{D}_S = \bb{0}$; flat: exact $\bb{D}_S$ & --- \\
    \bottomrule
  \end{tabular}
\end{table}

The fluid mesh follows the interface.
Its interior moves by a displacement Laplacian with quadratic inverse-distance diffusivity, and the inlet and outlet points are fixed (full variant) or move with the exact linear stretching $y \mapsto y\,h(t)/H$ (flat variant).
In the full variant this moves the cells next to the interface almost vertically.
They therefore become non-orthogonal to the tilting interface: $|\bb{k}|/(\hat{\bb{n}}\cdot\bb{d})$ reaches $0.24$ in the interface cells, about $13^\circ$, where $\bb{d}$ is the vector from the cell centre to the face centre and $\bb{k} = (\bb{I} - \hat{\bb{n}}\hat{\bb{n}})\cdot\bb{d}$.
The flat variant removes this at the interface; its inlet and outlet cells are slightly non-orthogonal ($|\bb{k}|/(\hat{\bb{n}}\cdot\bb{d}) \approx 0.15$).
The initial fields are zero.

\paragraph{Error measures.}
The errors are taken against the exact fields at the cell centres of the current mesh (fluid) and the initial mesh (solid).
They are reported as volume-weighted $L_2$ norms, $\|e\|_2 = (\sum_i V_i e_i^2/\sum_i V_i)^{1/2}$, and as maxima.
On the interface, the norms of the solid displacement and of the traction applied to the solid are weighted by the face areas of the initial solid interface; the applied traction is compared with $\bb{\sigma}_S\cdot\hat{\bb{n}}_S$, evaluated with the exact normal.
The meshes are one cell thick, with $n \times n/2$ fluid and $n \times n/4$ solid square cells for $n$ faces on the interface.
In a temporal study on a fixed mesh, successive time-step solutions are differenced as root-mean-square values over the cells; the fluid meshes of different time steps differ by the mesh motion, so the exact solution is first subtracted at each run's own cell centres.

\paragraph{Checks of the construction.}
The fields are built in sympy and differentiated in symengine.
Eleven properties are checked symbolically, or at $40$ significant digits at random points:
\begin{itemize}
  \item the inverse map, $\Xi(x^\Gamma(X,t),t) = X$, and $\bb{\nabla}\cdot(\bb{\nabla}\times\psi) = 0$ for any $\psi$, both symbolically;
  \item $h(x^\Gamma) = y^\Gamma$ and $\bb{\nabla}\cdot\bb{v}_F = 0$;
  \item the kinematic condition and no slip on $y = 0$;
  \item traction continuity, with the normal from Nanson's formula;
  \item zero solid interface shear for $m_1 = 0$, and the coefficient $c_1$;
  \item the clamps and the zero initial state.
\end{itemize}
The largest residuals are $3\times10^{-38}$ for the kinematic condition and $1.4\times10^{-38}$ for the traction (both variants), against an interface traction of order one.

The fields are evaluated in double precision by generated C++ (common-subexpression elimination).
A separate program checks the generated code independently of the symbolic pipeline:
\begin{itemize}
  \item fourth-order finite differences of the generated velocity and pressure reproduce $\bb{b}_F$ to $5\times10^{-10}$ and give $\bb{\nabla}\cdot\bb{v}_F$ to $4\times10^{-11}$;
  \item an independent implementation of the neo-Hookean law, applied to finite differences of $\bb{D}_S$, reproduces $\bb{\sigma}_S$ to $4\times10^{-10}$ and $\rho_{S,0}\bb{b}_S$ to $2\times10^{-7}$ (nested finite differences), against components of up to $8$;
  \item the kinematic condition, the interface normal and traction continuity hold, the last to $7\times10^{-15}$;
  \item at four reference points the generated values agree with the $40$-digit symbolic values to $2\times10^{-14}$, measured as $|a - b|/(1 + |b|)$.
\end{itemize}
All checks pass for both variants, and for a set with $\rho_F = 2$, $\rho_{S,0} = 3$ that tests the density scaling.
